% ============================================================================
%  Parâmetros do estudo — vêm ANTES da fundamentação teórica, porque todos os
%  métodos (ideal, natural e topo-plano) são demonstrados com estes valores.
% ============================================================================

\section{Parâmetros do Estudo}

\begin{frame}[allowframebreaks]{Parâmetros fixos}
    \small

    \textbf{Sinal original e espectro:}
    \begin{align*}
        x(t) &= \cos(2\pi \cdot 2t) + 0,5\cos(2\pi \cdot 18t) \\
        &\Updownarrow \\
        X(f) &= \frac{1}{2}[\delta(f-2) + \delta(f+2)]
        + \frac{0,5}{2}[\delta(f-18) + \delta(f+18)]
    \end{align*}

    \textbf{Filtro Anti-aliasing ($H_{AA}$):}
    \begin{align*}
        X_f(f) &= X(f) \cdot H_{AA}(f)
        \longrightarrow
        H_{AA}(f) =
        \begin{cases}
            1, & |f| \le 8\,\text{Hz}\\
            0, & |f| > 8\,\text{Hz}
        \end{cases}
        \quad \leftarrow f_c \\[2pt]
        X_f(f) &= \frac{1}{2}[\delta(f-2) + \delta(f+2)]
        \longleftrightarrow
        x_f(t) = \cos(2\pi \cdot 2t)
    \end{align*}

    \textbf{Parâmetros de Amostragem:}
    \[
        f_{max} = 2\,\text{Hz}
        \implies
        f_s > 2f_{max} = 4\,\text{Hz}
    \]

    \[
        f_s = 5\,\text{Hz}
        \implies
        T_s = \frac{1}{5}\,\text{s}
        \implies
        \text{banda de guarda} = 5 - 4 = 1\,\text{Hz}
    \]

    A partir dessa preparação, utiliza-se a equação filtrada em banda para realizar as amostragens: \begin{align*} x_f(t) = \cos(2\pi \cdot 2t) \end{align*}

\end{frame}

\begin{frame}{Parâmetros fixos}{Resumo}
    \small
    Estes são os valores usados em \textbf{todas} as demonstrações e figuras a seguir:
    \begin{center}
    \begin{tabular}{ll}
        \hline
        Sinal amostrado (pós anti-aliasing) & $x_f(t) = \cos(2\pi \cdot 2t)$ \\
        Frequência do sinal & $f_0 = f_{max} = 2$~Hz \\
        Corte do anti-aliasing & $f_c = 8$~Hz \\
        Frequência de amostragem & $f_s = 5$~Hz \\
        Período de amostragem & $T_s = 0{,}2$~s \\
        Banda de guarda & $f_s - 2f_{max} = 1$~Hz \\
        \hline
    \end{tabular}
    \end{center}

    \vspace{0.5em}
    Para cada método, o roteiro é o mesmo: descrição analítica geral, aplicação a
    $x_f(t)$ com estes parâmetros e, por fim, o filtro de reconstrução.
\end{frame}
